% !TEX TS-program = pdflatex
% !TeX program = pdflatex
% !TEX encoding = UTF-8
% !TEX spellcheck = en_US

\documentclass[xcolor=table]{beamer}


%\usepackage{fullpage}
%\usepackage[left=2.8cm,right=2.2cm,top=2 cm,bottom=2 cm]{geometry}
\setbeamersize{text margin left=10pt,text margin right=10pt}



%\usepackage[T1]{fontenc}
\usepackage{textcomp}
\usepackage[utf8]{inputenc}
%\usepackage[french,english]{babel}
\usepackage{arabtex}
\usepackage{txfonts}
\usepackage{tipa}
\usepackage[]{graphicx}
\usepackage{multirow}
\usepackage{hyperref}
\usepackage{colortbl}
\usepackage{tabularray}
\usepackage{listingsutf8}
\usepackage{wrapfig}
\usepackage{multicol}
\usepackage[export]{adjustbox} %for images in table, also for frame
\usepackage[many]{tcolorbox}
\usepackage{wasysym}
\usepackage[lined]{algorithm2e}
\usepackage{alltt} %verbatim with commands
\usepackage{longtable}
\usepackage{tabu}

\usepackage{CJKutf8}

\usepackage{amsmath,amssymb, amsfonts} 

%\usepackage{newtxtext,newtxmath}


%\usepackage{acolor}


\definecolor{lightblue}{HTML}{D0D2FF}
\definecolor{lightyellow}{HTML}{FFFFAA}
\definecolor{darkblue}{HTML}{0000BB}
\definecolor{olivegreen}{HTML}{006600}
\definecolor{darkgreen}{HTML}{008B45} %009B55
\definecolor{violet}{HTML}{6600CC}
\definecolor{deeppink}{HTML}{FF1493}
\definecolor{orangey}{HTML}{FFBB00}


\makeatletter
\def\beamer@calltheme#1#2#3{%
	\def\beamer@themelist{#2}
	\@for\beamer@themename:=\beamer@themelist\do
	{\IfFileExists{\beamer@themelocation/#3\beamer@themename.sty}{%
	\usepackage[{#1}]{\beamer@themelocation/#3\beamer@themename}}{\usepackage[{#1}]{#3\beamer@themename}}
		
	}%
}

\def\usefolder#1{
	\def\beamer@themelocation{#1}
}
\def\beamer@themelocation{}
\makeatother

\usefolder{../../extra/beamer-themes}

\usetheme{Karim} % Antibes Boadilla Warsaw Karim

\beamertemplatenavigationsymbolsempty

\let\oldcite\cite
\renewcommand{\cite}[1]{{\bfseries\color{blue}\oldcite{#1}}}

\tcbuselibrary{listings}


%\renewcommand{\baselinestretch}{1.5}

\def\supit#1{\raisebox{0.8ex}{\small\it #1}\hspace{0.05em}}

\AtBeginSection{%
	\begin{frame}
		\sectionpage
	\end{frame}
}

\newcommand{\rottext}[2]{%
	\rotatebox{90}{%
	\begin{minipage}{#1}%
		\raggedleft#2%
	\end{minipage}%
	}%
}


\institute{ %
Laboratoire de la Communication dans les Systèmes Informatiques (LCSI)
\par
École  nationale Supérieure d'Informatique (ESI), Algiers, Algeria
}
\author[Abdelkrime Aries (2023/2024)] %
{Abdelkrime Aries}
%\titlegraphic{\includegraphics[height=1cm]{../img/esi-logo.png}%\hspace*{4.75cm}~


\date{Academic year: 2023/2024} %\today

\titlegraphic{%
	\includegraphics[height=1cm]{../../img/misc/esi-logo.png}%
	\hspace{2cm}%
	\includegraphics[height=1cm]{../../img/misc/lcsi-logo.png}%
	\hspace{2cm}%
	\includegraphics[height=1.5cm]{../../img/misc/esi.nlp-logo.png}%
}

%\setbeamertemplate{headline}{}

\newcommand{\kurl}[1]{{\scriptsize\bfseries\color{orangey}\url{#1}}}

\newcommand{\keyword}[1]{\textcolor{red}{\bfseries\itshape #1}}
\newcommand{\expword}[1]{\textcolor{olivegreen}{#1}}
\newcommand{\optword}[1]{\textcolor{violet}{\bfseries #1}}

\makeatletter
\newcommand\mysphere{%
	\parbox[t]{10pt}{\raisebox{0.2pt}{\beamer@usesphere{item projected}{bigsphere}}}}
\makeatother

%\let\oldtabular\tabular
%\let\endoldtabular\endtabular
%\renewenvironment{tabular}{\rowcolors{2}{white}{lightblue}\oldtabular\rowcolor{blue}}{\endoldtabular}


%\NoAutoSpacing %french autospacing after ":"

\def\graphpath{}

\newcommand{\changegraphpath}[1]{\def\graphpath{#1}}


\newcommand{\vgraphpage}[2][.84\textheight]{%
%	\begin{center}%
		\includegraphics[height=#1]{\graphpath #2}%
%	\end{center}%
}

\newcommand{\hgraphpage}[2][\textwidth]{%
%	\begin{center}%
		\includegraphics[width=#1]{\graphpath #2}%
%	\end{center}%
}

\newcommand{\graphpage}[2][]{%
	\includegraphics[#1]{\graphpath #2}%
}

\bibliographystyle{apalike}

\newcommand{\insertbibliography}[2]{
	\appendix
	\section*{Bibliography}
	\nocite{#2}
%	\makeatletter % to change template
%	\setbeamertemplate{headline}[default] % not mandatory, but I though it was better to set it blank
%	\def\beamer@entrycode{\vspace*{-\headheight}} % here is the part we are interested in :)
%	\makeatother
	\begin{multicols*}{2}[\frametitle{\insertsection} \usebeamertemplate*{frametitle}]%\usebeamertemplate*{frametitle}\frametitle{Références}
		\tiny
		\bibliography{#1}
	\end{multicols*}
}

\definecolor{my-grey}{RGB}{233, 233, 233}

\newcommand{\insertlicence}{
	\begin{frame}[plain]
	\frametitle{License: CC-BY 4.0}

	\begin{tcolorbox}[colback=cyan,
		colframe=cyan,  
		arc=0pt,outer arc=0pt,
		valign=top, 
		halign=center,
		width=\textwidth]
		
		\includegraphics[width=.5cm]{../../img/licence/cc_icon_white_x2.png}
		\includegraphics[width=.5cm]{../../img/licence/attribution_icon_white_x2.png}
		
		\color{white}
		\bfseries Attribution 4.0 International (CC BY 4.0) \\
		\tiny \url{https://creativecommons.org/licenses/by/4.0/deed.en}
		
	\end{tcolorbox}\vspace{-.5cm}
	\begin{tcolorbox}[colback=my-grey,
		colframe=my-grey,  
		center, arc=0pt,outer arc=0pt,
		valign=top, 
		halign=left,
		width=\textwidth]
		
		\tiny
		
		\begin{center}
			\bfseries\Large
			You are free to:
		\end{center}
		
		\begin{minipage}{0.83\textwidth}
			\begin{itemize}
				\item[] \textbf{Share} — copy and redistribute the material in any medium or format
				\item[] \textbf{Adapt} — remix, transform, and build upon the material
				for any purpose, even commercially.
			\end{itemize}
		
			\vspace{12pt}
		
			\textit{The licensor cannot revoke these freedoms as long as you follow the license terms.}
		\end{minipage}
		\begin{minipage}{0.15\textwidth}
			\includegraphics[width=\textwidth]{../../img/licence/FreeCulturalWorks_seal_x2.jpg}
		\end{minipage}
	
		
		\begin{center}
			\bfseries\Large
			Under the following terms:
		\end{center}
		
		\begin{itemize}
			\item[] \textbf{Attribution} — You must give appropriate credit, provide a link to the license, and indicate if changes were made. You may do so in any reasonable manner, but not in any way that suggests the licensor endorses you or your use.
			\item[] \textbf{No additional restrictions} — You may not apply legal terms or technological measures that legally restrict others from doing anything the license permits.
		\end{itemize}
		
	\end{tcolorbox}
	
%	\begin{center}
%		\bfseries Attribution 4.0 International (CC BY 4.0)
%		\url{https://creativecommons.org/licenses/by/4.0/deed.fr}
%	\end{center}

%	\tiny
%
%	Vous êtes autorisé à : 
%	\begin{itemize}
%		\item \textbf{Partager} — copier, distribuer et communiquer le matériel par tous moyens et sous tous formats
%		\item \textbf{Adapter} — remixer, transformer et créer à partir du matériel
%	\end{itemize}
%	
%	Selon les conditions suivantes : 
%	\begin{itemize}
%		\item \textbf{Attribution} — Vous devez créditer l'Œuvre, intégrer un lien vers la licence et indiquer si des modifications ont été effectuées à l'Oeuvre. Vous devez indiquer ces informations par tous les moyens raisonnables, sans toutefois suggérer que l'Offrant vous soutient ou soutient la façon dont vous avez utilisé son Oeuvre.
%		\item \textbf{Pas d'Utilisation Commerciale} — Vous n'êtes pas autorisé à faire un usage commercial de cette Oeuvre, tout ou partie du matériel la composant. 
%		\item \textbf{Pas de restrictions complémentaires} — Vous n'êtes pas autorisé à appliquer des conditions légales ou des mesures techniques qui restreindraient légalement autrui à utiliser l'Oeuvre dans les conditions décrites par la licence.
%	\end{itemize}

	\end{frame}
}

\settowidth{\leftmargini}{\usebeamertemplate{itemize item}}
\addtolength{\leftmargini}{\labelsep}

\AtBeginDocument{
	\newcolumntype{L}[2]{>{\vbox to #2\bgroup\vfill\flushleft}p{#1}<{\egroup}} 
	
	\begin{frame}[plain]
		\maketitle
	\end{frame}

	\insertlicence
}


% needs etoolbox; to break links after -
\appto\UrlBreaks{\do\-}


\makeatletter
\newcommand{\xRightarrow}[2][]{\ext@arrow 0359\Rightarrowfill@{#1}{#2}}
\makeatother


\usefonttheme{structurebold}
%\usefonttheme{professionalfonts}


% figure caption
\setlength\abovecaptionskip{2pt}

% equations
\AtBeginDocument{
\abovedisplayskip=0pt
\abovedisplayshortskip=0pt
\belowdisplayskip=0pt
\belowdisplayshortskip=0pt
}


\title[ESI - NLP: 04- Language models]%
{Natural Language Processing\\Chapter 04\\Language models} 

\changegraphpath{../img/lm/}

\begin{document}
	
\begin{frame}
\frametitle{Natural Language Processing}
\framesubtitle{Language models: Introduction}

\begin{itemize}
	\item \optword{Sentence probability}: $\color{red} P(S) = P(w_1, w_2, ..., w_n) $
	\begin{itemize}
		\item Machine translation: \\
		\expword{My tall brother \textrightarrow\ P(Mon grand frère) \textgreater\ P(Mon haut frère)}
		\item Spell checking: \\
		\expword{P(They say something.) \textgreater\ P(They says something.)}
		\item Speech recognition: \\
		\expword{P(I say) \textgreater\ P(Eye say)}
	\end{itemize}
	\item \optword{word occurrence probability}: $\color{red} P(w_i | w_1, \ldots, w_{i-1}) $
	\begin{itemize}
		\item Auto-completion: \\
		\expword{P(automatic information system) \textgreater\ P(automatic information processing)}
		\item Automatic text generation
	\end{itemize}
\end{itemize}

\end{frame}


\begin{frame}
\frametitle{Natural Language Processing}
\framesubtitle{Language models: Plan}

\begin{multicols}{2}
%	\small
\tableofcontents
\end{multicols}
\end{frame}

%===================================================================================
\section{N-gram model}
%===================================================================================

\begin{frame}
\frametitle{Language models}
\framesubtitle{N-gram model}

\begin{block}{Compound probability formula}
	\[ P(w_1 \ldots w_m) =  P(w_1) P(w_2 | w_1) P(w_3 | w_1, w_2) \ldots P(w_n | w_1, \ldots, w_{m-1}) \]
\end{block}


\begin{exampleblock}{Example of a sentence's probability}
	\[ P(\text{\textit{I work at ESI}}) =  P(I) P(work | I) P(\text{\textit{at}} | \text{\textit{I work}}) P(ESI | \text{\textit{I work at}}) \]
	
	\[ P(ESI | \text{\textit{I work at}}) = \frac{C(\text{\textit{I work at ESI}})}{C( \text{\textit{I work at}})} \] 
	
	where 
	
	$C$: returns the frequency of an expression in a training corpus
\end{exampleblock}

\begin{itemize}
	\item Many possible sentences
	\item A big corpus is needed to estimate this probability (infinity of possible sentences)
\end{itemize}

\end{frame}

\subsection{Formulation}

\begin{frame}
\frametitle{Language models: N-gram}
\framesubtitle{Formulation: Markov property}

\begin{block}{Markov property}
	A present state depends only on the past one. 
	\[%
	P(w_i | w_1,\ldots, w_{i-1}) \approx P(w_i | w_{i-1})
	\]
	General case with $n-1$ past states
	\[%
	P(w_i | w_1,\ldots, w_{i-1}) \approx P(w_i | w_{i-n+1}, \ldots, w_{i-1})
	\]
\end{block}

\begin{block}{Proability estimation using N-grams}
	\[
	P(w_1,\ldots, w_{m}) \approx \prod_{i=1}^m P(w_i | w_{i-k+1}, \ldots, w_{i-1})
	\]
\end{block}

\end{frame}

\begin{frame}
\frametitle{Language models: N-gram}
\framesubtitle{Formulation: Some models}

\begin{itemize}
	\item \optword{Uni-gram model}: $P(w_i | w_1,\ldots, w_{i-1}) \approx P(w_i)$
	\item \optword{Bi-gram model}: $P(w_i | w_1,\ldots, w_{i-1}) \approx P(w_i | w_{i-1})$
	\item \optword{Tri-gram model}:  $P(w_i | w_1,\ldots, w_{i-1}) \approx P(w_i | w_{i-2}, w_{i-1})$
	\item Google Books Ngram Viewer
	\begin{itemize}
		\item \url{https://books.google.com/ngrams}
		\item Character-based language models from books 
		\item Free download: \url{https://storage.googleapis.com/books/ngrams/books/datasetsv3.html}
	\end{itemize}
\end{itemize}

\end{frame}

\begin{frame}
\frametitle{Language models: N-gram}
\framesubtitle{Formulation: Google Books Ngram Viewer (Humor)}

\begin{center}
	\vgraphpage{humor/humor-ngram.png}
\end{center}

\end{frame}

\begin{frame}
\frametitle{Language models: N-gram}
\framesubtitle{Formulation: Estimation}

\begin{itemize}
	\item Using a training corpus with enough data
	\item The beginning and ending of each sentence is marked by \keyword{\textless s\textgreater} and \keyword{\textless/s\textgreater} respectively (one time of bi-grams, 2 times for tri-grams, etc.)
	\item \keyword{Maximum likelihood estimator (MLE)}
\end{itemize}

\begin{block}{Proability estimation using MLE}
	{\small \[%
	P(w_i | w_{i-n+1},\ldots, w_{i-1}) = \frac{C(w_{i-n+1} \ldots w_{i-1} w_i)}{\sum_i C(w_{i-n+1} \ldots w_{i-1} w_i)}
	= \frac{C(w_{i-n+1} \ldots w_{i-1} w_i)}{C(w_{i-n+1} \ldots w_{i-1})}
	\]}
	Where $C$ is the frequency of N-grams in the training corpus
	\[%
	\text{Bi-grams: } P(w_i | w_{i-1}) = \frac{C(w_{i-1} w_i)}{C(w_{i-1})}
	\]
\end{block}

\end{frame}

\begin{frame}
\frametitle{Language models: N-gram}
\framesubtitle{Formulation: Example}

\begin{exampleblock}{Example of a training corpus (Bi-grams)}
	\begin{itemize}
		\item \textless s\textgreater a computer can help you \textless/s\textgreater
		\item \textless s\textgreater he wants to help you \textless/s\textgreater
		\item \textless s\textgreater he wants a computer \textless/s\textgreater
		\item \textless s\textgreater he can swim \textless/s\textgreater
	\end{itemize}
\end{exampleblock}

\begin{itemize}
	\item $P(can | he) = \frac{C(he\ can)}{C(he)} = \frac{1}{3}$
	\item $P(\text{\textit{\textless s\textgreater he can help you \textless/s\textgreater}}) = 
	\underbrace{P(he|\text{\textit{\textless s\textgreater}})}_{\frac{3}{4}}
	\underbrace{P(can|he)}_{\frac{1}{3}} 
	\underbrace{P(help|can)}_{\frac{1}{2}} 
	\underbrace{P(you|help)}_{\frac{2}{2}}
	\underbrace{P(\text{\textit{\textless/s\textgreater}}|you)}_{\frac{2}{2}} = 
	\frac{1}{8}
	$
\end{itemize}


\end{frame}

\subsection{Smoothing}

\begin{frame}
\frametitle{Language models: N-gram}
\framesubtitle{Smoothing}

\begin{itemize}
	
	\item Using little Ns $ \Longrightarrow $ Information loss
	\begin{itemize}
		\item Languages allow long term dependencies
		\item \expword{\underline{The computer} which I used yesterday at ESI during the course session \underline{has crashed}}
	\end{itemize}

	\item Using big Ns $ \Longrightarrow $ High model complexity
	\begin{itemize}
		\item This needs a bigger corpus
		\item N-grams representation: $V^N$ where $V$ is the vocabulary size and $N$ is the number of grams
	\end{itemize}

	\item Problem of missing N-grams in the training corpus
	\begin{itemize}
		\item $P(\text{\textit{\textless s\textgreater he wants to swim \textless/s\textgreater}}) = 
		P(he|\text{\textit{\textless s\textgreater}}) P(wants|he) P(to|wants) P(swim|to)  P(\text{\textit{\textless/s\textgreater}}|swim) = $
		
		$\frac{3}{4} \frac{2}{3} \frac{1}{2} \frac{0}{1} \frac{1}{1} = 0$
	\end{itemize}
	\item The intuition is to borrow a small portion of the existing N-gram probabilities to form the missing N-gram probability
\end{itemize}

\end{frame}

\begin{frame}
\frametitle{Language models: N-gram}
\framesubtitle{Smoothing: Lidstone}

\begin{block}{Lidstone smoothing (Bi-grams as example)}
	\[%
	P(w_i | w_{i-1}) = \frac{C(w_{i-1} w_i) + \alpha}{C(w_{i-1}) + \alpha V}
	\]
	Where $V$ is the vocabulary size of the model
	
	$\alpha = 1$ : \keyword{Laplace smoothing} 
	
	$\alpha = 0.5$ : \keyword{Jeffreys-Perks law}
\end{block}

\begin{exampleblock}{Example: Laplace smoothing}
	\begin{itemize}
		\item The corpus contains 9 different words
		\item $|V| = 9 + 2 = 11$ (Counting start and end markers)
		\item $P(\text{\textit{\textless s\textgreater he wants to swim \textless/s\textgreater}}) = 
		P(he|\text{\textit{\textless s\textgreater}}) P(wants|he) P(to|wants) P(swim|to) P(\text{\textit{\textless/s\textgreater}}|swim) = 
		\frac{3 + 1}{4 + 11} \frac{2 + 1}{3 + 11} \frac{1 + 1}{2 + 11} \frac{0 + 1}{1 + 11} \frac{1 + 1}{1 + 11}
		= \frac{4}{15} \frac{3}{14} \frac{2}{13} \frac{1}{12} \frac{2}{12} $
	\end{itemize}
\end{exampleblock}

\end{frame}

\begin{frame}
\frametitle{Language models: N-gram}
\framesubtitle{Smoothing: Interpolation}

\begin{block}{Interpolation (Tri-grams as example)}
	\[%
	P_{I}(w_i | w_{i-2} w_{i-1}) = 
	\lambda_3 P(w_i | w_{i-2} w_{i-1}) 
	+ \lambda_2 P(w_i | w_{i-1}) 
	+ \lambda_1 P(w_i) 
	\]
	
	Where $\sum_j \lambda_j = 1$
	
	$\lambda_3$, $\lambda_2$ and $\lambda_1$ are estimated using another tuning corpus
\end{block}

\begin{exampleblock}{Example: Bi-gram interpolation}
	\begin{itemize}
		\item Lets say $\lambda_2=0.8$ and $\lambda_1=0.2$
		\item P(\text{\textit{\textless/s\textgreater}}) will be calculated; we will consider "\textless/s\textgreater" in uni-grams
		\item $P(\text{\textit{\textless s\textgreater he wants to swim \textless/s\textgreater}}) = 
		P(he|\text{\textit{\textless s\textgreater}}) P(wants|he) P(to|wants) P(swim|to) P(\text{\textit{\textless/s\textgreater}}|swim) = 
		(0.8\frac{3}{4} + 0.2\frac{3}{21}) (0.8\frac{2}{3} + 0.2\frac{2}{21}) (0.8\frac{1}{2} + 0.2\frac{1}{21}) (0.8\frac{0}{1} + 0.2\frac{1}{21}) (0.8\frac{1}{1} + 0.2\frac{4}{21})$
	\end{itemize}
\end{exampleblock}

\end{frame}

\begin{frame}
\frametitle{Language models: N-gram}
\framesubtitle{Smoothing: Katz back-off}

\begin{block}{Katz back-off (Tri-grams as example)}
	\[%
	P_{BO}(w_i | w_{i-2} w_{i-1}) = 
	\begin{cases}
	P^*(w_i | w_{i-2} w_{i-1}) & \text{if } C(w_{i-2} w_{i-1} w_i) > 0 \\
	\alpha(w_{i-2} w_{i-1}) P_{BO}(w_i | w_{i-1}) & \text{otherwise}
	\end{cases}
	\]
	
	Where: 
	
	$P^*$ is the reduced probability (the reduction will be distributed over the lower order N-gram probabilities)
	
	$\alpha$ is a function which distributes the reduction according to the context
\end{block}

%\begin{exampleblock}{Example: Bi-gram interpolation}
%	\begin{itemize}
%		\item Lets say $\lambda_2=0.8$ and $\lambda_1=0.2$
%		\item P(\text{\textit{\textless/s\textgreater}}) will be calculated; we will consider "\textless/s\textgreater" in uni-grams
%		\item $P(\text{\textit{\textless s\textgreater he wants to swim \textless/s\textgreater}}) = 
%		P(he|\text{\textit{\textless s\textgreater}}) P(wants|he) P(to|wants) P(swim|to) P(\text{\textit{\textless/s\textgreater}}|swim) = 
%		(0.8\frac{3}{4} + 0.2\frac{3}{21}) (0.8\frac{2}{3} + 0.2\frac{2}{21}) (0.8\frac{1}{2} + 0.2\frac{1}{21}) (0.8\frac{0}{1} + 0.2\frac{1}{21}) (0.8\frac{1}{1} + 0.2\frac{4}{21})$
%	\end{itemize}
%\end{exampleblock}

\end{frame}


%===================================================================================
\section{Neural models}
%===================================================================================

%\begin{frame}
%\frametitle{Langage models}
%\framesubtitle{Neural models}
%
%\begin{itemize}
%	\item 
%\end{itemize}
%
%\end{frame}

\subsection{Multi-layers perceptron}

\begin{frame}
\frametitle{Language models: Neural models}
\framesubtitle{Multi-layers perceptron: Study case}

\begin{itemize}
	\item \cite{2003-bengio-al} model will be shown as a study case
	\item First, n-grams number $n$ must be chosen; in this case, the number of input words is $n-1$
	\item The vector $w$ is the concatenation of $n-1$ past words' embeddings
	\item The embedding size $d$ must be chosen  
	\item Words are encoded using One-Hot format (each vector has a size of the vocabulary $V$)
\end{itemize}

\begin{block}{MLP-based language model}
	$
	P(h_i|h_{i-n+1},\ldots, h_{i-1}) = 
	Softmax \left(
	(b + w A) 
	+ 
	W\ Tanh(u + w T)
	\right)
	$
	
	Where $b \in \mathbb{R}^{V},\, A \in \mathbb{R}^{(n-1) \times d \times V},\, u \in \mathbb{R}^{H},\, T \in \mathbb{R}^{(n-1) \times d \times H},\, W \in \mathbb{R}^{V \times H}$
\end{block}

\end{frame}

\begin{frame}
\frametitle{Language models: Neural models}
\framesubtitle{Multi-layers perceptron: Model}

\begin{center}
	\vgraphpage{mlp-model.pdf}
\end{center}

\end{frame}

\subsection{Recurrent neural networks}

\begin{frame}
\frametitle{Language models: Neural models}
\framesubtitle{Recurrent neural networks: Study case}

\begin{itemize}
	\item \cite{2010-mokolov-al} model will be shown as a study case
	\item For each instant $t$, a state $s_t$ is calculated based on the previous state $s_{t-1}$ and the currant word $w_t$
	\item The state $s_t$ is used to estimate the probabilities $y_t$ of each vocabulary word
\end{itemize}

\begin{block}{RNN-based language model}
	$x_t = s_{t-1} \bullet w_t$
	
	$s_t = \sigma(x_t W)$
	
	$y_t = softmax(s_t U)$
	
	Where $w_t \in \mathbb{R}^{V},\, s_t \in \mathbb{R}^{H},\, W \in \mathbb{R}^{(H+V)\times H},\, U \in \mathbb{R}^{H\times V}$
\end{block}

\end{frame}

\begin{frame}
\frametitle{Language models: Neural models}
\framesubtitle{Recurrent neural networks: Model}

\begin{center}
	\vgraphpage{rnn-model.pdf}
\end{center}

\end{frame}

\subsection{Some improvements}

\begin{frame}
\frametitle{Language models: Neural models}
\framesubtitle{Some improvements}

\begin{itemize}
	\item Limited context size in MLP-based models
	\begin{itemize}
		\item Use RNNs
	\end{itemize}
	\item Vanishing gradient problem in RNNs
	\begin{itemize}
		\item Use advanced RNNs: \keyword{LSTM} and \keyword{GRU}
		\item Limit context size
	\end{itemize}
	\item Out-of-vocabulary words
	\begin{itemize}
		\item Limit vocabulary size  et mark infrequent words as \optword{\textlangle UNK\textrangle}
	\end{itemize}
\end{itemize}

\end{frame}

%===================================================================================
\section{Evaluation}
%===================================================================================

\subsection{Approaches}

\begin{frame}
\frametitle{Language models: Evaluation}
\framesubtitle{Approaches}

\begin{itemize}
	\item \optword{Extrinsic Evaluation}
	\begin{itemize}
		\item Evaluate the model against another task: its effect on the task
		\item Example, \expword{Machine translation quality using a language model}
		\item Very expensive evaluation
	\end{itemize}
	\item \optword{Intrinsic Rating}
	\begin{itemize}
		\item Evaluate the model against its capacity to represent the language
		\item Example, \expword{Comparing two models based on their ability to represent a test dataset}
		\item Does not guarantee good model performance for a given task
	\end{itemize}
\end{itemize}

\end{frame}

\subsection{Perplexity}

\begin{frame}
\frametitle{Language models: Evaluation}
\framesubtitle{Perplexity}

\begin{itemize}
	\item Measuring the prediction quality of a model on a test corpus
	\item Using the estimated probability on a test corpus of size $N$
	\item Model with minimal perplexity is best
	\item Must include the end of one sentence and the beginning of the next in model training (since perplexity treats the whole corpus as a single string)
\end{itemize}

\begin{block}{Perplexity}
	\begin{center}
		$PP(w) = \sqrt[N]{\frac{1}{P(w_1 w_2 \ldots w_N)}}$
		
		$PP(w) = \sqrt[N]{\prod\limits_{i=1}^{N}\frac{1}{P(w_i | w_1 \ldots w_{i-1})}}$
	\end{center}
\end{block}

\end{frame}


\begin{frame}
\frametitle{Language models}
\framesubtitle{Some humor}

\begin{center}
	\vgraphpage{humor/humor-lm.png}
\end{center}

\end{frame}

\insertbibliography{NLP04}{*}

\end{document}

